% GENERATED FILE. Edit canonical lessons, then run npm run generate:books.
% canonical-source-hash: fnv1a64:817bac2e08ae337e
% canonical-lessons: RU-C218-raznyy, RU-C218-odinakovyy, RU-C218-pokhozhiy, RU-C218-glavnyy, RU-C218-obychnyy

\chapter{Different, Same, Similar, Main, Ordinary}
\label{ch:ru-different-same-similar-main-ordinary}
\hlchaptermodality{218}

\begin{chapteropening}
\textbf{By the end of this chapter:} \emph{I can say different, the same, similar, main and ordinary.}

Complete the last lesson of chapter 218: I can say different, the same, similar, main and ordinary.
\end{chapteropening}

\section[\texorpdfstring{\ru{разный} (ráznyy)}{разный (ráznyy)}]{\ru{разный} (ráznyy) — different, various}

\label{lesson:RU-C218-raznyy}



\begin{quote}
Before the new one: say the Russian for fresh, then the Russian for modern.
\end{quote}

\subsection*{You'll want to know: \ru{разный}}
\textbf{\ru{разный}} (\emph{ráznyy}) — \textquotedblleft{}different, various\textquotedblright{}.

\begin{grammarlens}[title={What you've built}]
One more word for this chapter.
\end{grammarlens}

\subsection*{Your turn}
\begin{itemize}
\raggedright
  \item \emph{Say it:} \emph{ráznyy}
  \item \emph{Say it:} \emph{ráznyy}, once more
  \item \emph{Say it:} say \emph{ráznyy}
  \item \emph{Recall:} say the Russian for fresh, then the Russian for modern, then say \emph{ráznyy} again
\end{itemize}

\begin{tcolorbox}[breakable,colback=teal!4,colframe=teal!35!black,title={Before you move on}]
What does \ru{разный} mean? (\textquotedblleft{}Different, various\textquotedblright{}.) Say it once more.
\end{tcolorbox}
\section[\texorpdfstring{\ru{одинаковый} (odinákovyy)}{одинаковый (odinákovyy)}]{\ru{одинаковый} (odinákovyy) — the same, identical}

\label{lesson:RU-C218-odinakovyy}



\begin{quote}
Before the new one: say the Russian for modern, then the Russian for different.
\end{quote}

\subsection*{You'll want to know: \ru{одинаковый}}
\textbf{\ru{одинаковый}} (\emph{odinákovyy}) — \textquotedblleft{}the same, identical\textquotedblright{}.

\begin{grammarlens}[title={What you've built}]
One more word for this chapter.
\end{grammarlens}

\subsection*{Your turn}
\begin{itemize}
\raggedright
  \item \emph{Say it:} \emph{odinákovyy}
  \item \emph{Say it:} \emph{odinákovyy}, once more
  \item \emph{Say it:} say \emph{odinákovyy}
  \item \emph{Recall:} say the Russian for modern, then the Russian for different, then say \emph{odinákovyy} again
\end{itemize}

\begin{tcolorbox}[breakable,colback=teal!4,colframe=teal!35!black,title={Before you move on}]
What does \ru{одинаковый} mean? (\textquotedblleft{}The same, identical\textquotedblright{}.) Say it once more.
\end{tcolorbox}
\section[\texorpdfstring{\ru{похожий} (pokhózhiy)}{похожий (pokhózhiy)}]{\ru{похожий} (pokhózhiy) — similar, alike}

\label{lesson:RU-C218-pokhozhiy}



\begin{quote}
Before the new one: say the Russian for different, then the Russian for the same.
\end{quote}

\subsection*{You'll want to know: \ru{похожий}}
\textbf{\ru{похожий}} (\emph{pokhózhiy}) — \textquotedblleft{}similar, alike\textquotedblright{}.

\begin{grammarlens}[title={What you've built}]
One more word for this chapter.
\end{grammarlens}

\subsection*{Your turn}
\begin{itemize}
\raggedright
  \item \emph{Say it:} \emph{pokhózhiy}
  \item \emph{Say it:} \emph{pokhózhiy}, once more
  \item \emph{Say it:} say \emph{pokhózhiy}
  \item \emph{Recall:} say the Russian for different, then the Russian for the same, then say \emph{pokhózhiy} again
\end{itemize}

\begin{tcolorbox}[breakable,colback=teal!4,colframe=teal!35!black,title={Before you move on}]
What does \ru{похожий} mean? (\textquotedblleft{}Similar, alike\textquotedblright{}.) Say it once more.
\end{tcolorbox}
\section[\texorpdfstring{\ru{главный} (glávnyy)}{главный (glávnyy)}]{\ru{главный} (glávnyy) — main}

\label{lesson:RU-C218-glavnyy}



\begin{quote}
Before the new one: say the Russian for the same, then the Russian for similar.
\end{quote}

\subsection*{You'll want to know: \ru{главный}}
\textbf{\ru{главный}} (\emph{glávnyy}) — \textquotedblleft{}main\textquotedblright{}.

\begin{grammarlens}[title={What you've built}]
One more word for this chapter.
\end{grammarlens}

\subsection*{Your turn}
\begin{itemize}
\raggedright
  \item \emph{Say it:} \emph{glávnyy}
  \item \emph{Say it:} \emph{glávnyy}, once more
  \item \emph{Say it:} say \emph{glávnyy}
  \item \emph{Recall:} say the Russian for the same, then the Russian for similar, then say \emph{glávnyy} again
\end{itemize}

\begin{tcolorbox}[breakable,colback=teal!4,colframe=teal!35!black,title={Before you move on}]
What does \ru{главный} mean? (\textquotedblleft{}Main\textquotedblright{}.) Say it once more.
\end{tcolorbox}
\section[\texorpdfstring{\ru{обычный} (obýchnyy)}{обычный (obýchnyy)}]{\ru{обычный} (obýchnyy) — ordinary, usual}

\label{lesson:RU-C218-obychnyy}



\begin{quote}
Before the new one: say the Russian for similar, then the Russian for main.
\end{quote}

\subsection*{You'll want to know: \ru{обычный}}
\textbf{\ru{обычный}} (\emph{obýchnyy}) — \textquotedblleft{}ordinary, usual\textquotedblright{}.

\begin{grammarlens}[title={What you've built}]
One more word for this chapter.
\end{grammarlens}

\subsection*{Your turn}
\begin{itemize}
\raggedright
  \item \emph{Say it:} \emph{obýchnyy}
  \item \emph{Say it:} \emph{obýchnyy}, once more
  \item \emph{Say it:} say \emph{obýchnyy}
  \item \emph{Recall:} say the Russian for similar, then the Russian for main, then say \emph{obýchnyy} again
\end{itemize}

\begin{tcolorbox}[breakable,colback=teal!4,colframe=teal!35!black,title={Before you move on}]
What does \ru{обычный} mean? (\textquotedblleft{}Ordinary, usual\textquotedblright{}.) Say it once more.
\end{tcolorbox}
